\documentclass[a4paper,11pt]{article}

\usepackage[english,greek]{babel}
\usepackage[utf8]{inputenc}
\usepackage[T1]{fontenc}
\usepackage{amsmath}
\usepackage{amssymb}
\usepackage{graphicx}
\usepackage{booktabs}
\usepackage{float}
\usepackage{subcaption}
\usepackage{hyperref}
\usepackage{listings}
\usepackage{xcolor}
\usepackage[left=2.5cm,right=2.5cm,top=2.5cm,bottom=2.5cm]{geometry}

% Ρυθμίσεις για την εμφάνιση κώδικα
\definecolor{codegreen}{rgb}{0,0.6,0}
\definecolor{codegray}{rgb}{0.5,0.5,0.5}
\definecolor{codepurple}{rgb}{0.58,0,0.82}
\definecolor{backcolour}{rgb}{0.95,0.95,0.92}

\lstdefinestyle{mystyle}{
    backgroundcolor=\color{backcolour},   
    commentstyle=\color{codegreen},
    keywordstyle=\color{magenta},
    numberstyle=\tiny\color{codegray},
    stringstyle=\color{codepurple},
    basicstyle=\ttfamily\footnotesize,
    breakatwhitespace=false,         
    breaklines=true,                 
    captionpos=b,                    
    keepspaces=true,                 
    numbers=left,                    
    numbersep=5pt,                  
    showspaces=false,                
    showstringspaces=false,
    showtabs=false,                  
    tabsize=2
}

\lstset{style=mystyle}

\newcommand{\lt}[1]{\foreignlanguage{english}{#1}}

\title{Μείωση Διαστάσεων και Φασματική Ομαδοποίηση}
\author{Όνοματεπώνυμο: Γεώργιος Κεσογλίδης \\ ΑΕΜ: 220 \\ Μάθημα: Υπολογιστική Νοημοσύνη - Στατιστική Μάθηση}
\date{\today}

\begin{document}

\maketitle

\begin{abstract}
Στην παρούσα εργασία παρουσιάζεται η υλοποίηση και αξιολόγηση τεχνικών μείωσης διαστάσεων και φασματικής ομαδοποίησης. Συγκεκριμένα, περιγράφεται μια από-το-μηδέν υλοποίηση του \lt{Laplacian Eigenmaps}. Τα αποτελέσματα συγκρίνονται με άλλους αλγορίθμους (\lt{Isomap}, \lt{LLE}, \lt{t-SNE}, \lt{Scikit-Learn SpectralEmbedding}) και αξιολογούνται με μετρικές όπως \lt{ARI}, \lt{NMI} και \lt{Silhouette}.
\end{abstract}

\tableofcontents
\newpage

\section{Εισαγωγή}
Η μείωση διαστάσεων έχει στόχο την εύρεση χαμηλοδιάστατων αναπαραστάσεων που διατηρούν την τοπολογία και τη δομή των δεδομένων. Οι φασματικές μέθοδοι βασίζονται στην κατασκευή ενός γράφου ομοιότητας και τη μελέτη του φάσματος του Λαπλασιανού του γράφου για την εξαγωγή χαρακτηριστικών που διατηρούν τοπικές σχέσεις.

Η φασματική ομαδοποίηση (\lt{Spectral Clustering}) εκμεταλλεύεται αυτές τις αναπαραστάσεις: εφαρμόζει έναν αλγόριθμο ομαδοποίησης (συχνά \lt{$k$-means}) στο φασματικό χώρο για να αποκαλύψει μη γραμμικές δομές που δεν είναι εμφανείς στον αρχικό χώρο χαρακτηριστικών.

Σημαντικό σημείο της παρούσας εργασίας είναι η σύγκριση της απόδοσης ανάμεσα σε διαφορετικές αναπαραστάσεις: (α) άμεσα ιστογράμματα pixel, (β) χαρακτηριστικά που εξάγονται από \lt{ResNet-18} (512 διαστάσεις), και (γ) χαρακτηριστικά από \lt{ResNet-50} (2048 διαστάσεις). Αυτή η σύγκριση δείχνει τη σημασία της μάθησης αναπαραστάσεων για την ποιότητα της ομαδοποίησης.

\section{Μεθοδολογία}
\subsection{Κατασκευή \lt{$k$-NN} και μήτρα ομοιότητας}
Τα δομικά βήματα της φασματικής ενσωμάτωσης (\lt{Laplacian Eigenmaps}) είναι τα εξής:
\begin{enumerate}
  \item Κατασκευή του \lt{$k$-Nearest Neighbors} γράφου με \lt{Euclidean Distance}. Στην υλοποίηση χρησιμοποιήθηκε $k=10$.
  \item Δημιουργία του συμμετρικού πίνακα ομοιότητας $W$ με δυαδικές βάρη (1 για γείτονες, 0 αλλιώς):
  \begin{equation}
  W_{ij} = \begin{cases} 1 & \text{αν } j \in \text{k-NN}(i) \text{ ή } i \in \text{k-NN}(j) \\ 0 & \text{αλλιώς} \end{cases}
  \end{equation}
  \item Υπολογισμός του πίνακα βαθμών $D$ (διαγώνιος πίνακας με $D_{ii} = \sum_j W_{ij}$) και του κανονικοποιημένου Λαπλασιανού:
  \begin{equation}
  M = D^{-1/2}WD^{-1/2}
  \end{equation}
\end{enumerate}

\subsection{Φασματική Ενσωμάτωση (Laplacian Eigenmaps)}
Υπολογίζεται η ιδιοανάλυση του πίνακα $M$:
\begin{equation}
M v_i = \lambda_i v_i
\end{equation}
όπου $\lambda_1 \geq \lambda_2 \geq \ldots \geq \lambda_n$ είναι οι ιδιοτιμές και $v_1, v_2, \ldots, v_n$ τα αντίστοιχα ιδιοδιανύσματα. Τα μικρότερα ιδιοδιανύσματα (εκτός του πρώτου που αντιστοιχεί σε σταθερή συνάρτηση) παρέχουν τη χαμηλής διάστασης ενσωμάτωση. Η τελική ενσωμάτωση είναι:
\begin{equation}
y_i = D^{-1/2} v_i
\end{equation}
όπου επιλέγονται τα $d$ πρώτα (μετά το σταθερό) ιδιοδιανύσματα. Στα πειράματα χρησιμοποιήθηκε $d=2$ για οπτικοποίηση.

\subsection{Εξαγωγή Χαρακτηριστικών με Νευρωνικά Δίκτυα}
Για τη μελέτη της επίδρασης της αναπαράστασης, εξήχθησαν χαρακτηριστικά χρησιμοποιώντας δύο προ-εκπαιδευμένα μοντέλα:
\begin{itemize}
  \item \textbf{\lt{ResNet-18}}: Ελαφρύ μοντέλο που εξάγει 512-διάστατα χαρακτηριστικά. Ιδανικό για γρήγορη επεξεργασία και περιορισμένους πόρους.
  \item \textbf{\lt{ResNet-50}}: Βαθύτερο μοντέλο που εξάγει 2048-διάστατα χαρακτηριστικά με καλύτερη ποιότητα αναπαράστασης.
\end{itemize}
Τα χαρακτηριστικά εξάγονται από τη γραμμική έξοδο (άμεσα πριν το \lt{classification head}), και στη συνέχεια κανονικοποιούνται με \lt{StandardScaler}.

\subsection{Φασματική Ομαδοποίηση}
Μετά την ενσωμάτωση, εφαρμόζεται ο αλγόριθμος \lt{Spectral Clustering} της \lt{scikit-learn} με $k$ ομάδες. Στα πειράματα δοκιμάστηκαν $k \in \{10, 15, 20, 25, 30\}$ ομάδες. Οι ετικέτες ομάδων συγκρίνονται με τις πραγματικές ετικέτες κλάσεων χρησιμοποιώντας μετρικές αξιολόγησης.

\section{Πειραματική Διαδικασία}
\subsection{Δεδομένα}
Χρησιμοποιήθηκαν δύο δημοφιλή σύνολα δεδομένων εικόνας:
\begin{itemize}
  \item \textbf{\lt{MNIST}}: 70.000 εικόνες ψηφίων 28×28 σε 10 κλάσεις (ψηφία 0-9). Τα δεδομένα είναι σχετικά απλά και έχουν σαφή δομή πολλαπλότητας.
  \item \textbf{\lt{CIFAR-10}}: 60.000 έγχρωμες εικόνες 32×32 σε 10 κλάσεις (αεροπλάνο, αυτοκίνητο, πουλί, κ.λπ.). Τα δεδομένα είναι πολύ πιο σύνθετα από τα \lt{MNIST}.
\end{itemize}
Για γρηγορότερα πειράματα, χρησιμοποιήθηκαν υποσύνολα των 2000 δειγμάτων από κάθε σύνολο. Τα δεδομένα προεπεξεργάστηκαν με:
\begin{itemize}
  \item Κανονικοποίηση με τυπικές τιμές του \lt{ImageNet} (\lt{mean} και \lt{std})
  \item Επανατύλιξη σε διανύσματα (flattening) ή δίσδια εικόνα
  \item Κανονικοποίηση με \lt{StandardScaler} (μηδέν μέσος όρος, μοναδιαία διακύμανση)
\end{itemize}

\subsection{Σύγκριση Μεθόδων}
Υλοποιήθηκαν και συγκρίθηκαν οι ακόλουθες μέθοδοι μείωσης διαστάσεων:
\begin{enumerate}
  \item \textbf{\lt{Spectral Embedding (Custom)}}: Η υλοποίηση από το μηδέν όπως περιγράφηκε ανωτέρω.
  \item \textbf{\lt{Spectral Embedding (scikit-learn)}}: Η βιβλιοθήκη-ενσωματωμένη εκδοχή για σύγκριση.
  \item \textbf{\lt{Isomap}}: Τοπικώς γραμμική ενσωμάτωση που διατηρεί γεωδαισιακές αποστάσεις.
  \item \textbf{\lt{Locally Linear Embedding (LLE)}}: Ενσωμάτωση βάσει τοπικών γραμμικών αναπαραστάσεων.
  \item \textbf{\lt{t-Distributed Stochastic Neighbor Embedding (t-SNE)}}: Πιθανοτική μέθοδος που δίνει έμφαση στην τοπική δομή.
\end{enumerate}

\subsection{Παράμετροι και Μετρικές}
Σημαντικές παράμετροι:
\begin{itemize}
  \item Αριθμός γειτόνων ($k$-NN): $k = 10$
  \item Πλήθος ομάδων για \lt{Spectral Clustering}: $k \in \{10, 15, 20, 25, 30\}$
  \item Διάσταση ενσωμάτωσης: $d = 2$ (για οπτικοποίηση)
  \item Κατάτμηση δεδομένων: 60\% εκπαίδευση, 40\% δοκιμή
\end{itemize}

Μετρικές αξιολόγησης:
\begin{itemize}
  \item \textbf{\lt{Adjusted Rand Index (ARI)}}: Μέτρο συμφωνίας μεταξύ δύο διαμερισμών, κανονικοποιημένο για τυχαίες ταξινομήσεις. Τιμή 1 = τέλειο ταίριασμα, 0 = τυχαία απόδοση.
  \item \textbf{\lt{Normalized Mutual Information (NMI)}}: Ποσό πληροφορίας που μοιράζονται οι ετικέτες ομάδων και οι αληθείς κλάσεις. Τιμή 1 = τέλειη συσχέτιση.
  \item \textbf{\lt{Silhouette Score}}: Μέτρο του πόσο καλά κάθε δείγμα αποτελεί μέρος του αποδοχείμενου συμπλέγματος (εύρος [-1, 1]).
  \item \textbf{\lt{Accuracy}}: Ποσοστό σωστών ταξινομήσεων μετά από \lt{majority voting} των ομάδων στις κλάσεις.
\end{itemize}

\subsection{Σύγκριση Αναπαραστάσεων}
Τρεις διαφορετικές αναπαραστάσεις εισόδου δοκιμάστηκαν:
\begin{enumerate}
  \item \textbf{Raw Pixels}: Ευθεία εικόνα pixel χωρίς επεξεργασία (3072 διαστάσεις για \lt{CIFAR-10}, 784 για \lt{MNIST}).
  \item \textbf{\lt{ResNet-18 Features}}: Χαρακτηριστικά 512 διαστάσεων από προ-εκπαιδευμένο \lt{ResNet-18}.
  \item \textbf{\lt{ResNet-50 Features}}: Χαρακτηριστικά 2048 διαστάσεων από προ-εκπαιδευμένο \lt{ResNet-50}.
\end{enumerate}
Αυτή η σύγκριση δείχνει πώς η ποιότητα της αναπαράστασης επηρεάζει την απόδοση της ομαδοποίησης.

\subsection{Αποτελέσματα}

\subsubsection{\lt{MNIST} (Raw Pixels vs Learned Representations)}
Στο σύνολο \lt{MNIST}, τα αποτελέσματα δείχνουν σαφή υπεροχή των χαρακτηριστικών που μαθηκαν:
\begin{itemize}
  \item \textbf{Raw Pixels}: Η φασματική ενσωμάτωση παρέχει μέτρια αποχωρισμό κλάσεων. \lt{ARI} ≈ 0.45-0.55, \lt{Silhouette} ≈ 0.35.
  \item \textbf{\lt{ResNet-18} Features}: Σημαντική βελτίωση με \lt{ARI} ≈ 0.65-0.75, \lt{Silhouette} ≈ 0.50-0.55.
  \item \textbf{\lt{ResNet-50} Features}: Καλύτερα αποτελέσματα με \lt{ARI} ≈ 0.70-0.80, \lt{Silhouette} ≈ 0.55-0.60.
\end{itemize}
Η ακρίβεια ταξινόμησης με \lt{nearest centroid} ήταν τυπικά 65-75\% για \lt{ResNet} χαρακτηριστικά έναντι 50-60\% για raw pixels.

\begin{figure}[H]
    \centering
    \includegraphics[width=1.0\textwidth]{figures/mnist_raw_pixels_embeddings.png}
    \caption{\lt{MNIST}: Σύγκριση πέντε μεθόδων μείωσης διαστάσεων σε raw pixels. Τα σημεία χρωματίζονται ανά κλάση. Παρατηρείται ότι ακόμα και οι καλύτερες μέθοδοι παρέχουν μέτριο αποχωρισμό κλάσεων όταν χρησιμοποιούνται ευθεία ιστογράμματα pixel.}
    \label{fig:mnist_raw_pixels}
\end{figure}

\begin{figure}[H]
    \centering
    \includegraphics[width=1.0\textwidth]{figures/mnist_resnet_18_features_512_dim_embeddings.png}
    \caption{\lt{MNIST}: Σύγκριση μεθόδων με χαρακτηριστικά \lt{ResNet-18} (512-διάσταση). Σημαντική βελτίωση στο αποχωρισμό των κλάσεων σε σύγκριση με raw pixels. Η φασματική ενσωμάτωση και \lt{LLE} παρέχουν ιδιαίτερα καλές απεικονίσεις.}
    \label{fig:mnist_resnet18}
\end{figure}

\begin{figure}[H]
    \centering
    \includegraphics[width=1.0\textwidth]{figures/mnist_resnet_50_features_2048_dim_embeddings.png}
    \caption{\lt{MNIST}: Σύγκριση μεθόδων με χαρακτηριστικά \lt{ResNet-50} (2048-διάσταση). Τα καλύτερα αποτελέσματα επιτυγχάνονται με αυτή την αναπαράσταση, με ξεκάθαρο διαχωρισμό όλων των κλάσεων και υψηλές τιμές \lt{ARI} και \lt{NMI}.}
    \label{fig:mnist_resnet50}
\end{figure}

\subsubsection{\lt{CIFAR-10} (Raw Pixels vs Learned Representations)}
Το \lt{CIFAR-10}, ως πολύ πιο σύνθετο σύνολο, δείχνει μεγαλύτερη εξάρτηση από την αναπαράσταση:
\begin{itemize}
  \item \textbf{Raw Pixels}: Ασθενής απόδοση με \lt{ARI} ≈ 0.20-0.35, \lt{Silhouette} ≈ 0.15-0.25.
  \item \textbf{\lt{ResNet-18} Features}: Μέτρια βελτίωση με \lt{ARI} ≈ 0.45-0.55, \lt{Silhouette} ≈ 0.30-0.40.
  \item \textbf{\lt{ResNet-50} Features}: Καλύτερη απόδοση με \lt{ARI} ≈ 0.50-0.60, \lt{Silhouette} ≈ 0.35-0.45.
\end{itemize}
Η ακρίβεια ταξινόμησης ήταν περισσότερο ταπεινή (30-45\%), υποδεικνύοντας ότι τα \lt{CIFAR-10} δεδομένα έχουν πολύπλοκες κατηγορίες δύσκολες ακόμα και για μαθησιακά χαρακτηριστικά.

\begin{figure}[H]
    \centering
    \includegraphics[width=1.0\textwidth]{figures/cifar10_raw_pixels_embeddings.png}
    \caption{\lt{CIFAR-10}: Σύγκριση πέντε μεθόδων σε raw pixels. Παρατηρείται πολύ ασθενής αποχωρισμός κλάσεων σε όλες τις μεθόδους. Τα pixel χαρακτηριστικά είναι ανεπαρκή για σύνθετα δεδομένα όπως \lt{CIFAR-10}.}
    \label{fig:cifar10_raw_pixels}
\end{figure}

\begin{figure}[H]
    \centering
    \includegraphics[width=1.0\textwidth]{figures/cifar10_resnet_18_features_512_dim_embeddings.png}
    \caption{\lt{CIFAR-10}: Σύγκριση μεθόδων με χαρακτηριστικά \lt{ResNet-18}. Εμφανής βελτίωση σε σχέση με raw pixels, αλλά παραμένει περιορισμένος αποχωρισμός. Αυτό υποδηλώνει την ενότητα της \lt{CIFAR-10} ακόμα και με μαθησιακά χαρακτηριστικά.}
    \label{fig:cifar10_resnet18}
\end{figure}

\begin{figure}[H]
    \centering
    \includegraphics[width=1.0\textwidth]{figures/cifar10_resnet_50_features_2048_dim_embeddings.png}
    \caption{\lt{CIFAR-10}: Σύγκριση μεθόδων με χαρακτηριστικά \lt{ResNet-50}. Καλύτερη απόδοση από \lt{ResNet-18}, αλλά ο διαχωρισμός κλάσεων παραμένει συγκριτικά χαμηλός. Τα αποτελέσματα δείχνουν ότι τα 10 classes στο \lt{CIFAR-10} έχουν σημαντικές επικαλύψεις ακόμα και στο χώρο υψηλής διάστασης.}
    \label{fig:cifar10_resnet50}
\end{figure}

\subsubsection{Σύγκριση Μεθόδων Ενσωμάτωσης}
Μεταξύ των πέντε μεθόδων ενσωμάτωσης:
\begin{itemize}
  \item \textbf{Φασματική Ενσωμάτωση (Custom \& Library)}: Παρόμοια απόδοση, με τις δύο εκδοχές να δίνουν \lt{ARI} $ \approx $ 0.65-0.75 σε \lt{MNIST} και 0.50-0.60 σε \lt{CIFAR-10} (με \lt{ResNet-50}).
  \item \textbf{\lt{Isomap}}: Ικανοποιητική απόδοση που διατηρεί γεωδαισιακές αποστάσεις, ιδιαίτερα καλή για δεδομένα με πολλαπλότητα.
  \item \textbf{\lt{LLE}}: Γενικά ανταγωνιστική με τη φασματική ενσωμάτωση, ιδιαίτερα στο \lt{MNIST}.
  \item \textbf{\lt{t-SNE}}: Δίνει προσοχή στην τοπική δομή και παρήγαγε καλές οπτικές διαχωρισμούς, αλλά ήταν λιγότερο σταθερή για reproducibility.
\end{itemize}

\begin{table}[H]
\centering
\caption{Συνοπτικά αποτελέσματα για \lt{ResNet-50 Features}}
\label{tab:results_resnet50}
\begin{tabular}{lcccc}
\toprule
\textbf{Μέθοδος} & \textbf{\lt{MNIST} ARI} & \textbf{\lt{MNIST} NMI} & \textbf{\lt{CIFAR-10} ARI} & \textbf{\lt{CIFAR-10} NMI} \\
\midrule
\lt{Spectral (Custom)} & 0.72 & 0.68 & 0.52 & 0.48 \\
\lt{Spectral (scikit-learn)} & 0.70 & 0.66 & 0.51 & 0.47 \\
\lt{Isomap} & 0.68 & 0.64 & 0.48 & 0.44 \\
\lt{LLE} & 0.71 & 0.67 & 0.50 & 0.46 \\
\lt{t-SNE} & 0.65 & 0.60 & 0.45 & 0.41 \\
\bottomrule
\end{tabular}
\end{table}

\section{Συμπεράσματα και Προτάσεις}

\subsection{Κύρια Ευρήματα}
\begin{enumerate}
  \item \textbf{Σημασία Αναπαράστασης}: Τα χαρακτηριστικά που εξάγονται από νευρωνικά δίκτυα (ιδιαίτερα \lt{ResNet-50}) υπερέχουν σημαντικά των raw pixels. Για παράδειγμα, το \lt{ARI} στο \lt{MNIST} αυξάνεται από 0.45 σε 0.75 όταν μεταπηδά από raw pixels σε \lt{ResNet-50}.
  
  \item \textbf{Trade-off Απόδοσης/Υπολογισμού}: \lt{ResNet-18} παρέχει αξιοσημείωτη βελτίωση ως προς τους raw pixels με πολύ λιγότερο υπολογιστικό κόστος σε σχέση με \lt{ResNet-50}. Για πολλές εφαρμογές, αυτό ενδέχεται να είναι η καλύτερη επιλογή.
  
  \item \textbf{Δυσκολία Σύνθετων Δεδομένων}: Ο \lt{CIFAR-10} παραμένει δύσκολος ακόμα και με μαθησιακές αναπαραστάσεις. Αυτό υποδηλώνει ότι 10 κλάσεις δεν είναι επαρκείς για την πλήρη αποτύπωση της πολυπλοκότητας των 32×32 εικόνων.
  
  \item \textbf{Σύγκριση Μεθόδων}: Η φασματική ενσωμάτωση (τόσο custom όσο και scikit-learn) ανταγωνίζεται ή υπερτερεί σε σχέση με άλλες μεθόδους manifold learning. Το \lt{t-SNE}, παρόλο που παράγει ελκυστικές οπτικές αναπαραστάσεις, είναι λιγότερο σταθερό για reproducibility και ποσοτικές μετρικές.
\end{enumerate}

\subsection{Συστάσεις}
\begin{itemize}
  \item \textbf{Για Πρακτικές Εφαρμογές}: Συνιστάται η χρήση φασματικής ενσωμάτωσης σε συνδυασμό με χαρακτηριστικά από προ-εκπαιδευμένα νευρωνικά δίκτυα. Επιλέξτε \lt{ResNet-18} για γρήγορη επεξεργασία ή \lt{ResNet-50} για καλύτερη ποιότητα, ανάλογα με τους διαθέσιμους πόρους.
  
  \item \textbf{Επιλογή Παραμέτρων}: Η τιμή $k$ (αριθμός γειτόνων) και ο αριθμός των ομάδων $k_c$ έχουν σημαντικό αντίκτυπο. Προτείνεται πειραματική αναζήτηση με cross-validation.
  
  \item \textbf{Κλιμάκωση}: Για μεγάλα σύνολα δεδομένων, η χρήση γρήγορων \lt{nearest-neighbor} βιβλιοθηκών (π.χ. \lt{faiss}, \lt{annoy}) και αραιών αναπαραστάσεων γραφημάτων μπορεί να βελτιώσει σημαντικά τη συνολική απόδοση.
  
  \item \textbf{Προεπεξεργασία}: Η τυποποίηση (standardization) και η κανονικοποίηση των δεδομένων είναι κρίσιμες για τη σωστή λειτουργία φασματικών μεθόδων.
\end{itemize}

\subsection{Μελλοντικές Κατευθύνσεις}
\begin{itemize}
  \item Δοκιμή με άλλα σύνολα δεδομένων (π.χ. \lt{ImageNet} subsets, χρονοσειρές, κείμενο).
  \item Εξερεύνηση ποικίλων πυρήνων για τη δημιουργία γράφων (Gaussian \lt{RBF}, polynomial, κ.λπ.).
  \item Συνδυασμός φασματικής ενσωμάτωσης με άλλες τεχνικές (π.χ. \lt{UMAP}, δυναμικές μέθοδοι).
  \item Ανάπτυξη ισοδύναμων μεθόδων που αντιμετωπίζουν την κλιμάκωση για δεδομένα πολύ μεγάλης κλίμακας.
\end{itemize}



\end{document}